\mathtgt
\section{The Full Target: \texorpdfstring{\FCdot}{FCdot}}\label{sec:target}

This section gives the whole of \FCdot. It adds four features to the fragment of Section~\ref{sec:miniature}, fields, recursive types, intersections, and self-bounds, and one more kind of template. Each subsection says what the feature adds and why. Figures~\ref{fig:target-syntax}, \ref{fig:evidence}, \ref{fig:terms}, and~\ref{fig:machine} give the complete calculus and can be read on their own.

\begin{figure}
\small
\color{tgtcol}
\[
\begin{array}{@{}lrcl@{}}
\text{Types} & S,T & ::= & \bot \mid \sel{x}{\ell} \mid \pit[x]{S}{T} \mid \obj[x]{\tel} \\
& \top & \mathrel{:=} & \obj{\tnil} \\
\text{Propositions} & P & ::= & \ple{S}{T} \mid \peq{S}{T} \mid \phas{\ell} \mid \pbnd{T} \\
\text{Telescopes} & \tel & ::= & \tnil \mid \tcons{\tel}{P} \\[2pt]
\text{Inclusion} & e & ::= & \ev{refl}\,T \mid \ev{trans}\,e\,e \mid \ev{top}\,T \mid \ev{bot}\,T \mid \ev{eqToLe}\,\varphi \mid \ev{pi}\,e\,e \\
& & \mid & \ev{obj}\,\tel\,m \mid \ev{pair}\,\tel\,\tel\,e\,e \mid \ev{member}\,a\,e\,i \mid \ev{bound}\,\tel\,i \mid \ev{intoBnd}\,e \\
\text{Equality} & \varphi & ::= & \ev{refl}\,T \mid \ev{symm}\,\varphi \mid \ev{trans}\,\varphi\,\varphi \mid \ev{def}\,x\,\ell \mid \ev{member}\,a\,e\,i \\
\text{Presence} & h & ::= & \ev{member}\,a\,e\,i \mid \ev{field}\,\ell \\
\text{Holes} & H & ::= & \ev{le}\,j \mid \ev{eq}\,j \mid \ev{eqSym}\,j \\
\text{Sides} & q & ::= & \ev{none} \mid \ev{some}\,e \\
\text{Morphisms} & m & ::= & [\,\tau, \ldots, \tau\,] \\
\text{Templates} & \tau & ::= & \ev{le}\,q\,H\,q' \mid \ev{leTrans}\,\tau\,\tau \mid \ev{eq}\,j\,b \mid \ev{has}\,j \mid \ev{bnd}\,e \\[2pt]
\text{Atoms} & a,b & ::= & x \mid \ev{cast}\,a\,e \mid \ev{fold}\,\tel\,a \mid \ev{unfold}\,a \mid \ev{both}\,\tel\,\tel\,a\,a \\
\text{Terms} & t,u & ::= & a \mid v \mid a\,b \mid \ev{proj}\,a\,\ell\,h \mid \klet x = t \kin u \mid \ev{cast}\,t\,e \\
\text{Values} & v & ::= & \lambda(x : S).\,t \mid \lit[x]{W}{F} \mid \ev{cast}\,v\,e \\
\text{Witnesses} & W & ::= & \tnil \mid W, \ell = T \\
\text{Fields} & F & ::= & \tnil \mid F, \ell = t \\
\text{Contexts} & \Gamma & ::= & \cdot \mid \Gamma, x : T \mid \Gamma, x : T \mathrel{:=} W ; \vec{\ell}
\end{array}
\]
\color{black}
\caption{Syntax of \FCdot. The binders are those of the fragment: $\pit[x]{S}{T}$ binds $x$ in $T$, $\obj[x]{\tel}$ binds the self $x$ in $\tel$, and $\lit[x]{W}{F}$ binds the self $x$ in $W$ and $F$. A telescope carried by evidence or an atom has the self $\self$. In $\ev{pi}\,e\,e'$ the second component lives under the parameter.}
\label{fig:target-syntax}
\end{figure}

Figure~\ref{fig:target-syntax} gives the syntax. As in the source, everything is scoped intrinsically, and the two calculi share signatures, variables, renamings, and labels. The source is built on the binder infrastructure of the target, so the translation is the identity on signatures. This is a convenience of the mechanization, and nothing in the design of evidence depends on it.

\subsection{Fields and presence}

Fields bring a second kind of member name. As in the source, $A$ and $B$ name type members and $a$ and $b$ name fields, and the two share one name space of labels. From here on $\ell$ ranges over all labels, and every binder $x$ has a name $\sel{x}{\ell}$ for each of them. For a field $a$ the name $\sel{x}{a}$ is a type as well, and this subsection says what it stands for.

A term member $a$ of an object needs two facts: that the member exists, and what its type is. \FCdot keeps them apart. Presence, $\phas{a}$, is a proposition of its own. The type of the field is the block name $\sel{x}{a}$, which exists for every label like every other name, and what is known about it is stated by the other propositions. A field declaration translates to
\[
\trans{\dfld{a}{T}} = \obj{[\,\phas{a},\ \ple{\sel{\self}{a}}{\weak{\trans{T}}}\,]} ,
\]
which says that the field exists and that its type is included in $\trans{T}$.

A literal $\lit[x]{W}{F}$ now has fields $F$ besides its witnesses $W$, and binds its self $x$ in both. Its precise type is $\obj[x]{\ofLit{W}{F}}$, with one equality per witness followed by one presence per field:
\[
\ofLit{W}{F} = [\,\peq{\sel{x}{\ell_1}}{W(\ell_1)}, \ldots, \peq{\sel{x}{\ell_k}}{W(\ell_k)},\ \phas{\ell'_1}, \ldots, \phas{\ell'_n}\,] ,
\]
where $\ell_1, \ldots, \ell_k$ are the labels of $W$, oldest first, $\ell'_1, \ldots, \ell'_n$ are the labels of $F$, newest first, and $W(\ell)$ is the newest witness at $\ell$. The translation gives each field label a witness too, its declared type, so the name $\sel{x}{a}$ of a stored literal is defined as the type of the field. The fields are typed under the self binder $x$, bound transparently, and the body of field $\ell$ has the block name $\sel{x}{\ell}$ as its type. A context binder records the field labels of a transparent binder, written $x : T \mathrel{:=} W ; \vec\ell$.

Presence evidence has two forms. $\ev{member}\,a\,e\,i$ eliminates a presence proposition at an atom, like the other elimination rules. $\ev{field}\,\ell$ reads presence off a transparent binder that records the field. A projection $\ev{proj}\,a\,\ell\,h$ carries the presence evidence $h$ and has the block name $\sel{\rt{a}}{\ell}$ as its type. Converting that name to the declared field type takes a further cast, which the translation reads off the declaration. There are two more templates. $\ev{has}\,j$ proves a target presence from the source presence at position $j$, and $\ev{eq}\,j\,b$ proves a target equality from the source equality at $j$, flipped when $b$ is true.

\subsection{Recursive types}

In the fragment the bounds of a declaration could not mention the self. A recursive type $\dmu{x}{T}$ of DOT binds its own self, and that self is identified with the self of the object type. The type of Example~\ref{ex:hide}, with the intersection of the next subsection, becomes
\[
\trans{\dmu{x}{\dtyp{A}{\bot}{\top} \wedge \dfld{a}{x.A}}} = \obj{[\,\ple{\bot}{\sel{\self}{A}},\ \ple{\sel{\self}{A}}{\top},\ \phas{a},\ \ple{\sel{\self}{a}}{\sel{\self}{A}}\,]} ,
\]
where the last proposition relates two names of the self's own block. A declaration nested in the body may mention the outer self, and then the nested object type renames its own:
\[
\begin{array}{@{}l@{}}
\trans{\dmu{x}{\dtyp{A}{\bot}{\top} \wedge \dtyp{B}{\bot}{\dtyp{C}{x.A}{x.A}}}} \\
\quad = \obj{[\,\ldots,\ \ple{\sel{\self}{B}}{\obj[\self']{[\,\ple{\sel{\self}{A}}{\sel{\self'}{C}},\ \ple{\sel{\self'}{C}}{\sel{\self}{A}}\,]}}\,]} .
\end{array}
\]
With one fixed name for every self, the inner propositions could not mention the outer one. Opening a recursive type at a variable $x$, which is what \rn{Rec-E} does, replaces the self by $x$. In \FCdot this is an atom, $\ev{unfold}\,a$, which instantiates the self of $a$'s object type at the root of $a$ and weakens the result under a fresh self that nothing mentions:
\[
\inferrule{\atomj{\Gamma}{a}{\obj{\tel}}}{\atomj{\Gamma}{\ev{unfold}\,a}{\obj{\weak{(\inst{\tel}{\rt{a}})}}}}
\qquad
\inferrule{\atomj{\Gamma}{a}{\obj{\weak{(\inst{\tel}{\rt{a}})}}}}{\atomj{\Gamma}{\ev{fold}\,\tel\,a}{\obj{\tel}}}
\]
The atom $\ev{fold}\,\tel\,a$ is the inverse, the atom form of \rn{Rec-I}. Both erase to nothing and keep the root. The unfolded type is an object type whose propositions talk about the root instead of the self. A recursive type therefore needs no subtyping rule of its own in the target either, as it does not in the source.

Folding and unfolding change the type of an atom without evidence, so the chain of casts of an atom, which composes the head forms of its casts, cannot be typed between the plain types any more. Section~\ref{sec:canonical} types it in a mode at the root, in which object types are read with the self opened at the root, so that folding and unfolding are invisible.

\subsection{Intersections}

An intersection of two object types is the concatenation of their telescopes:
\[
\trans{\src{S \wedge T}} = \obj{\mathit{tel}\,\src{S} \mathbin{+\!\!+} \mathit{tel}\,\src{T}} .
\]
An intersection is therefore never an operation on types, only on lists of propositions. The projections \rn{And}$_1$ and \rn{And}$_2$ are object coercions whose templates copy one half of the source. The source rule \rn{And}, which pairs two subtypings out of one type, becomes $\ev{pair}\,\tel_1\,\tel_2\,e\,f$: from $e : S \le \obj{\tel_1}$ and $f : S \le \obj{\tel_2}$ it concludes $S \le \obj{\tel_1 \mathbin{+\!\!+} \tel_2}$. The source rule \rn{And-I}, which pairs two typings of one variable, becomes the atom $\ev{both}\,\tel_1\,\tel_2\,a\,b$, which requires $a$ and $b$ to have the same root. The telescope annotations make the endpoints computable from the arguments.

\subsection{Self-bounds}

Not every type is a declaration. An intersection $x.A \wedge \dfld{a}{T}$ has an operand that is a selection, and a recursive type $\dmu{x}{x.A}$ has a body that is one. Such a type cannot be split into propositions about named members. It enters a telescope as one self-bound, $\pbnd{T}$, a fourth kind of proposition, which says that the object itself is included in $T$:
\[
\trans{\src{x.A \wedge \dfld{a}{T}}} = \obj{[\,\pbnd{\weak{\sel{x}{A}}},\ \phas{a},\ \ple{\sel{\self}{a}}{\weak{\trans{T}}}\,]} .
\]
A bound has one elimination, one introduction, and one template.
\begin{itemize}
\item The cast $\ev{bound}\,\tel\,i$ takes the object type $\obj{\tel}$ through its $i$-th bound, $\obj{\tel} \le T$, when that bound is closed, $\pbnd{\weak{T}}$.
\item The evidence $\ev{intoBnd}\,e$ turns $S \le T$ into $S \le \obj{[\,\pbnd{\weak{T}}\,]}$, the object type with that one bound.
\item The template $\ev{bnd}\,e$ proves a target bound, where $e$ is closed evidence out of the source object type as a whole. This is the one template that reads the whole source rather than one of its propositions. It still eliminates nothing through the self, because $e$ is closed.
\end{itemize}
A bound may mention the self. That is what the body of $\dmu{x}{x.A}$ becomes: $\obj{[\,\pbnd{\sel{\self}{A}}\,]}$, the object type whose one proposition says that the object is included in its own member $A$. A cast through a bound reads only a bound whose type is closed, and the rules that prove a bound prove only such bounds. A bound that mentions the self is used by unfolding an atom of its type first, which instantiates the self at the root and makes every bound of the telescope closed.

\subsection{Composition templates}

A target inclusion may also be proven by a composition template $\ev{leTrans}\,\tau_1\,\tau_2$, which chains two templates over the same source, $\tau_1$ for $\ple{S}{M}$ and $\tau_2$ for $\ple{M}{T}$. The middle type $M$ may mention the self, so a composition template can chain facts that the source asserts about its own members, such as $\ple{\sel{\self}{a}}{\sel{\self}{A}}$ and $\ple{\sel{\self}{A}}{\sel{\self}{B}}$. The translation of this report produces no composition templates. They are there for recursive subtyping, to which Section~\ref{sec:related} returns, and the metatheory covers them.

\subsection{Evidence}

\begin{figure}
\small
\color{tgtcol}
\begin{mathpar}
\inferrule{ }{\lej{\Gamma}{\ev{refl}\,T}{T}{T}}
\and
\inferrule{\lej{\Gamma}{e}{S}{M} \\ \lej{\Gamma}{f}{M}{T}}{\lej{\Gamma}{\ev{trans}\,e\,f}{S}{T}}
\and
\inferrule{ }{\lej{\Gamma}{\ev{top}\,T}{T}{\top}}
\and
\inferrule{ }{\lej{\Gamma}{\ev{bot}\,T}{\bot}{T}}
\and
\inferrule{\eqj{\Gamma}{\varphi}{S}{T}}{\lej{\Gamma}{\ev{eqToLe}\,\varphi}{S}{T}}
\and
\inferrule{\lej{\Gamma}{e}{S_2}{S_1} \\ \lej{\Gamma, x : S_2}{f}{T_1}{T_2}}{\lej{\Gamma}{\ev{pi}\,e\,f}{\pit{S_1}{T_1}}{\pit{S_2}{T_2}}}
\and
\inferrule{\morj{\Gamma}{m}{\tel}{\tel'}}{\lej{\Gamma}{\ev{obj}\,\tel\,m}{\obj{\tel}}{\obj{\tel'}}}
\and
\inferrule{\lej{\Gamma}{e}{S}{\obj{\tel_1}} \\ \lej{\Gamma}{f}{S}{\obj{\tel_2}}}{\lej{\Gamma}{\ev{pair}\,\tel_1\,\tel_2\,e\,f}{S}{\obj{\tel_1 \mathbin{+\!\!+} \tel_2}}}
\and
\inferrule{\atomj{\Gamma}{a}{S} \\ \lej{\Gamma}{e}{S}{\obj{\tel}} \\ \tat{\tel}{i}{\ple{S'}{T'}}}
  {\lej{\Gamma}{\ev{member}\,a\,e\,i}{\inst{S'}{\rt{a}}}{\inst{T'}{\rt{a}}}}
\and
\inferrule{\tat{\tel}{i}{\pbnd{\weak{T}}}}{\lej{\Gamma}{\ev{bound}\,\tel\,i}{\obj{\tel}}{T}}
\and
\inferrule{\lej{\Gamma}{e}{S}{T}}{\lej{\Gamma}{\ev{intoBnd}\,e}{S}{\obj{[\,\pbnd{\weak{T}}\,]}}}
\\
\inferrule{ }{\eqj{\Gamma}{\ev{refl}\,T}{T}{T}}
\and
\inferrule{\eqj{\Gamma}{\varphi}{S}{T}}{\eqj{\Gamma}{\ev{symm}\,\varphi}{T}{S}}
\and
\inferrule{\eqj{\Gamma}{\varphi}{S}{M} \\ \eqj{\Gamma}{\psi}{M}{T}}{\eqj{\Gamma}{\ev{trans}\,\varphi\,\psi}{S}{T}}
\and
\inferrule{\Gamma(x) \ni \ell \mathrel{:=} T}{\eqj{\Gamma}{\ev{def}\,x\,\ell}{\sel{x}{\ell}}{T}}
\and
\inferrule{\atomj{\Gamma}{a}{S} \\ \lej{\Gamma}{e}{S}{\obj{\tel}} \\ \tat{\tel}{i}{\peq{S'}{T'}}}
  {\eqj{\Gamma}{\ev{member}\,a\,e\,i}{\inst{S'}{\rt{a}}}{\inst{T'}{\rt{a}}}}
\and
\inferrule{\atomj{\Gamma}{a}{S} \\ \lej{\Gamma}{e}{S}{\obj{\tel}} \\ \tat{\tel}{i}{\phas{\ell}}}
  {\hasj{\Gamma}{\ev{member}\,a\,e\,i}{\rt{a}}{\ell}}
\and
\inferrule{\Gamma(x) \text{ has fields } \vec{\ell} \\ \ell \in \vec{\ell}}{\hasj{\Gamma}{\ev{field}\,\ell}{x}{\ell}}
\\
\inferrule{ }{\Gamma \vdash \ev{none} : X \le X}
\and
\inferrule{\lej{\Gamma}{e}{A}{B}}{\Gamma \vdash \ev{some}\,e : \weak{A} \le \weak{B}}
\and
\inferrule{\morj{\Gamma}{\tau_i}{\tel}{P_i} \text{ for each } i < n}{\morj{\Gamma}{[\,\tau_0, \ldots, \tau_{n-1}\,]}{\tel}{[\,P_0, \ldots, P_{n-1}\,]}}
\and
\inferrule{\holej{\tel}{H}{X}{Y} \\ \Gamma \vdash q : S \le X \\ \Gamma \vdash q' : Y \le T}
  {\morj{\Gamma}{\ev{le}\,q\,H\,q'}{\tel}{\ple{S}{T}}}
\and
\inferrule{\morj{\Gamma}{\tau_1}{\tel}{\ple{S}{M}} \\ \morj{\Gamma}{\tau_2}{\tel}{\ple{M}{T}}}
  {\morj{\Gamma}{\ev{leTrans}\,\tau_1\,\tau_2}{\tel}{\ple{S}{T}}}
\and
\inferrule{\tat{\tel}{j}{\peq{X}{Y}}}
  {\morj{\Gamma}{\ev{eq}\,j\,\mathit{false}}{\tel}{\peq{X}{Y}}}
\and
\inferrule{\tat{\tel}{j}{\peq{X}{Y}}}
  {\morj{\Gamma}{\ev{eq}\,j\,\mathit{true}}{\tel}{\peq{Y}{X}}}
\and
\inferrule{\tat{\tel}{j}{\phas{\ell}}}
  {\morj{\Gamma}{\ev{has}\,j}{\tel}{\phas{\ell}}}
\and
\inferrule{\lej{\Gamma}{e}{\obj{\tel}}{T}}
  {\morj{\Gamma}{\ev{bnd}\,e}{\tel}{\pbnd{\weak{T}}}}
\\
\inferrule{\tat{\tel}{j}{\ple{X}{Y}}}{\holej{\tel}{\ev{le}\,j}{X}{Y}}
\and
\inferrule{\tat{\tel}{j}{\peq{X}{Y}}}{\holej{\tel}{\ev{eq}\,j}{X}{Y}}
\and
\inferrule{\tat{\tel}{j}{\peq{Y}{X}}}{\holej{\tel}{\ev{eqSym}\,j}{X}{Y}}
\end{mathpar}
\color{black}
\caption{Evidence typing. A hole $H$ reads the proposition of the source telescope that it names as an inclusion, $\holej{\tel}{H}{X}{Y}$.}
\label{fig:evidence}
\end{figure}

Figure~\ref{fig:evidence} gives the typing of all evidence. There are three families: inclusions $\lej{\Gamma}{e}{S}{T}$, equalities $\eqj{\Gamma}{\varphi}{S}{T}$, and presence $\hasj{\Gamma}{h}{x}{\ell}$, which says that the object bound at $x$ has a field $\ell$. The elimination $\ev{member}$ has one rule per kind of proposition it eliminates. It eliminates inclusions, equalities, and presence, but not bounds, which are cast through instead.

\subsection{Atoms, terms, and values}\label{sec:atoms}

\begin{figure}
\small
\color{tgtcol}
\begin{mathpar}
\inferrule{ }{\atomj{\Gamma}{x}{\Gamma(x)}}
\and
\inferrule{\atomj{\Gamma}{a}{S} \\ \lej{\Gamma}{e}{S}{T}}{\atomj{\Gamma}{\ev{cast}\,a\,e}{T}}
\and
\inferrule{\atomj{\Gamma}{a}{\obj{\tel}}}{\atomj{\Gamma}{\ev{unfold}\,a}{\obj{\weak{(\inst{\tel}{\rt{a}})}}}}
\and
\inferrule{\atomj{\Gamma}{a}{\obj{\weak{(\inst{\tel}{\rt{a}})}}}}{\atomj{\Gamma}{\ev{fold}\,\tel\,a}{\obj{\tel}}}
\and
\inferrule{\atomj{\Gamma}{a}{\obj{\tel_1}} \\ \atomj{\Gamma}{b}{\obj{\tel_2}} \\ \rt{b} = \rt{a}}{\atomj{\Gamma}{\ev{both}\,\tel_1\,\tel_2\,a\,b}{\obj{\tel_1 \mathbin{+\!\!+} \tel_2}}}
\\
\inferrule{\atomj{\Gamma}{a}{T}}{\tmj{\Gamma}{a}{T}}
\and
\inferrule{\valj{\Gamma}{v}{T}}{\tmj{\Gamma}{v}{T}}
\and
\inferrule{\atomj{\Gamma}{a}{\pit{S}{T}} \\ \atomj{\Gamma}{b}{S}}{\tmj{\Gamma}{a\,b}{\inst{T}{\rt{b}}}}
\and
\inferrule{\atomj{\Gamma}{a}{S} \\ \hasj{\Gamma}{h}{\rt{a}}{\ell}}{\tmj{\Gamma}{\ev{proj}\,a\,\ell\,h}{\sel{\rt{a}}{\ell}}}
\and
\inferrule{\tmj{\Gamma}{t}{T} \\ \tmj{\Gamma, x : T}{u}{\weak{U}}}{\tmj{\Gamma}{\klet x = t \kin u}{U}}
\and
\inferrule{\tmj{\Gamma}{t}{S} \\ \lej{\Gamma}{e}{S}{T}}{\tmj{\Gamma}{\ev{cast}\,t\,e}{T}}
\\
\inferrule{\tmj{\Gamma, x : S}{t}{T}}{\valj{\Gamma}{\lambda(x : S).\,t}{\pit{S}{T}}}
\and
\inferrule{\fldj{\Gamma, x : \obj[x]{\ofLit{W}{F}} \mathrel{:=} W ; \dom F}{F}}{\valj{\Gamma}{\lit[x]{W}{F}}{\obj[x]{\ofLit{W}{F}}}}
\and
\inferrule{\valj{\Gamma}{v}{S} \\ \lej{\Gamma}{e}{S}{T}}{\valj{\Gamma}{\ev{cast}\,v\,e}{T}}
\and
\inferrule{ }{\fldj{\Gamma}{\tnil}}
\and
\inferrule{\fldj{\Gamma}{F} \\ \tmj{\Gamma}{t}{\sel{x}{\ell}}}{\fldj{\Gamma}{F, \ell = t}}
\end{mathpar}
\color{black}
\caption{Typing of atoms, terms, values, and fields. In the field rules $x$ is the self binder of the literal.}
\label{fig:terms}
\end{figure}

Figure~\ref{fig:terms} gives the typing of atoms, terms, values, and fields. An atom is a variable under wrappers that erase to nothing: casts, folding and unfolding, and pairing. Its root is the variable underneath:
\[
\begin{array}{@{}lcl@{}}
\rt{x} & = & x \\
\rt{\ev{cast}\,a\,e} & = & \rt{a} \\
\rt{\ev{fold}\,\tel\,a} & = & \rt{a} \\
\rt{\ev{unfold}\,a} & = & \rt{a} \\
\rt{\ev{both}\,\tel_1\,\tel_2\,a\,b} & = & \rt{a} .
\end{array}
\]
In the last case $\rt{b}$ is the same variable whenever the atom is typed.

\subsection{Stores and the machine}

\begin{figure}
\small
\color{tgtcol}
\[
\begin{array}{@{}lcll@{}}
\mst{\sigma}{K}{\klet x = t \kin u} &\step& \mst{\sigma}{K \triangleright \ev{let}\,u}{t} \\
\mst{\sigma}{K}{\ev{cast}\,t\,e} &\step& \mst{\sigma}{K \triangleright \ev{cast}\,e}{t} \\
\mst{\sigma}{K \triangleright \ev{cast}\,e}{v} &\step& \mst{\sigma}{K}{\ev{cast}\,v\,e} \\
\mst{\sigma}{K \triangleright \ev{cast}\,e}{a} &\step& \mst{\sigma}{K}{\ev{cast}\,a\,e} \\
\mst{\sigma}{K \triangleright \ev{let}\,u}{v} &\step& \mst{\sigma, \mathit{core}(v)}{\weak{K}}{\mathit{adjust}(u, v)} \\
\mst{\sigma}{K \triangleright \ev{let}\,u}{a} &\step& \mst{\sigma}{K}{\inst{u}{a}} \\
\mst{\sigma}{K}{x\,b} &\step& \mst{\sigma}{K}{\inst{t_0}{b}} \\
& & \quad \text{if } \sigma(x) = \lambda(z : S_0).\,t_0 \\
\mst{\sigma}{K}{a\,b} &\step& \mst{\sigma}{K}{\inst{t_0}{b}} \\
& & \quad \text{if } \sigma(\rt{a}) = \lambda(z : S_0).\,t_0,\ a \text{ not a variable},\ \chn{\sigma}{a}{n}{\form{id}} \\
\mst{\sigma}{K}{a\,b} &\step& \mst{\sigma}{K}{\ev{cast}\,(\inst{t_0}{\ev{cast}\,b\,d})\,(\inst{c}{b})} \\
& & \quad \text{if } \sigma(\rt{a}) = \lambda(z : S_0).\,t_0,\ a \text{ not a variable},\ \chn{\sigma}{a}{n}{\form{pi}\,d\,c} \\
\mst{\sigma}{K}{\ev{proj}\,a\,\ell\,h} &\step& \mst{\sigma}{K}{\inst{t}{\rt{a}}} \\
& & \quad \text{if } \sigma(\rt{a}) = \lit{W}{F} \text{ and } F(\ell) = t
\end{array}
\]
\color{black}
\caption{The store machine of \FCdot. $\mathit{core}$ and $\mathit{adjust}$ are defined in Section~\ref{sec:mini-machine}. $\chn{\sigma}{a}{n}{F}$ is the normalization of the chain of casts of $a$ with fuel $n$, defined in Section~\ref{sec:normalizer}.}
\label{fig:machine}
\end{figure}

A store holds one value per binder of its signature. Store typing $\storej{\sigma}{\Gamma}$ requires every entry to be a lambda or a literal, typed in the transparent context of the entries before it, with the context recording the literal's witnesses and field labels. The self of a stored literal becomes its location:
\begin{mathpar}
\color{tgtcol}
\inferrule{ }{\storej{\cdot}{\cdot}}
\and
\inferrule{\storej{\sigma}{\Gamma} \\ v \text{ has no casts} \\ \valj{\Gamma}{v}{T}}{\storej{\sigma, v}{\Gamma, x : T \mathrel{:=} \mathit{witnesses}(v) ; \mathit{fields}(v)}}
\end{mathpar}
Here $\mathit{witnesses}(\lit[x]{W}{F}) = W$ and $\mathit{fields}(\lit[x]{W}{F}) = \dom F$, and both are empty for a lambda. Continuations and states are typed by the rules of Section~\ref{sec:mini-machine}, and a state is final when its continuation is empty and its term is a value or an atom.

Figure~\ref{fig:machine} gives the steps. They are those of the fragment, plus projection, which reads the field off the literal at the root and instantiates the field's self binder by the root. Substitution of an atom for a variable, $\inst{u}{a}$, puts the atom wherever an atom is expected, in terms and inside the evidence that eliminates at an atom. It puts the root of the atom into types and into the evidence that names a binder, $\ev{def}\,x\,\ell$, since those only ever see variables.

\subsection{Erasure}

Erasure into the runtime of Figure~\ref{fig:runtime} maps an atom to its root, drops every cast and every piece of evidence, keeps the fields of a literal, and drops cast frames:
\[
\begin{array}{@{}lcl@{}}
\erase{a} & = & \run{\rt{a}} \\
\erase{a\,b} & = & \run{\rt{a}\,\rt{b}} \\
\erase{\ev{proj}\,a\,\ell\,h} & = & \run{\rt{a}.\ell} \\
\erase{\klet x = t \kin u} & = & \run{\klet x = \erase{t} \kin \erase{u}} \\
\erase{\ev{cast}\,t\,e} & = & \erase{t} \\
\erase{\lambda(x : S).\,t} & = & \run{\lambda x.\,\erase{t}} \\
\erase{\lit[x]{W}{F}} & = & \run{\nu(x.\,\erase{F})} \\
\erase{\ev{cast}\,v\,e} & = & \erase{v} \\
\erase{F, \ell = t} & = & \run{\erase{F}, \ell = \erase{t}} \\
\erase{K \triangleright \ev{let}\,u} & = & \run{\erase{K} \triangleright (x.\,\erase{u})} \\
\erase{K \triangleright \ev{cast}\,e} & = & \erase{K}
\end{array}
\]
The empty field list and the empty continuation erase to themselves, and stores and states erase pointwise. A step of the machine either moves a cast frame, in which case the erased state does not change, or erases to exactly one runtime step. Cast-frame steps are bounded by the nesting of casts in the state, so they cannot be taken forever.

\subsection{The checker}

Every constructor of evidence and of atoms determines its endpoints from its arguments and the context. Object coercions carry their source telescope, pairing and folding carry their target telescopes, a basic template is checked from its hole outwards, and a composition template checks that its two templates prove one inclusion each and meet at the same middle type. The checker is therefore a family of synthesizing functions, one per judgment, each returning the typing derivation it found. Soundness holds by construction, and completeness, that every derivation is found, is a theorem proven by induction on derivations. Its top-level form is
\[
\lean{checkTm}\ \Gamma\ t\ T = \mathit{true} \iff \tmj{\Gamma}{t}{T} ,
\]
and the examples of Section~\ref{sec:mechanization} are checked by running it inside the kernel of the proof assistant. A checker matters for a calculus of evidence in the way that the Core Lint pass matters for GHC: the output of a translation can be validated independently of the translation.
